\documentclass[11pt,a4paper]{ctexart}
\usepackage[top=25mm,bottom=25mm,left=24mm,right=24mm]{geometry}
\usepackage{graphicx,booktabs,threeparttable,makecell,multirow,tabularx,array,enumitem,amsmath,amssymb,xcolor,float,fancyhdr,lastpage}
\usepackage[section]{placeins}
\usepackage{etoolbox}
\pretocmd{\subsection}{\FloatBarrier}{}{}
\usepackage[font=small,labelfont=bf,labelsep=quad,skip=4pt]{caption}
\usepackage[colorlinks=true,linkcolor=black,citecolor=black,urlcolor=blue]{hyperref}
\setmainfont{Times New Roman}
\setsansfont{Arial}
\ctexset{section={format=\Large\heiti\raggedright},subsection={format=\large\heiti\raggedright},
         subsubsection={format=\normalsize\heiti\raggedright}}
\newcolumntype{L}{>{\raggedright\arraybackslash}X}
\newcommand{\hlai}[1]{\textbf{#1}}
\renewcommand{\arraystretch}{1.18}
\setlist{nosep,leftmargin=2em}
\renewcommand{\topfraction}{0.9}\renewcommand{\bottomfraction}{0.8}\renewcommand{\textfraction}{0.07}
\renewcommand{\floatpagefraction}{0.8}\setcounter{topnumber}{3}\setcounter{totalnumber}{5}
\setlength{\emergencystretch}{2em}
\pagestyle{fancy}\fancyhf{}
\lhead{\small C3 FPS AI 声卡脚步增强功能 · 客观评测报告}\rhead{\small 第 \thepage\ 页 / 共 \pageref{LastPage} 页}
\renewcommand{\headrulewidth}{0.4pt}
\newcommand{\figslot}[4]{%
\begin{figure}[htbp]\centering
\IfFileExists{figures/#1.pdf}{\includegraphics[width=#2]{figures/#1.pdf}}{%
\IfFileExists{figures/#1.png}{\includegraphics[width=#2]{figures/#1.png}}{%
\fbox{\parbox[c][4.2cm][c]{#2}{\centering\small\color{gray}图位：\texttt{\detokenize{#1}}\\[2pt]按《图表需求说明》制作后放入 figures 目录}}}}
\caption{#3}\label{fig:#4}\end{figure}}
\newcommand{\keybox}[1]{\par\noindent{\parbox{\linewidth}{#1}}\par}


\renewcommand{\FloatBarrier}{}
\definecolor{ai}{HTML}{1F4E79}\definecolor{gun}{HTML}{D62728}
\newcommand{\tagbox}[1]{\colorbox{ai!10}{\strut\heiti\color{ai}\ #1\ }}
\newcommand{\greybox}[1]{\par\noindent\colorbox{black!7}{\parbox{\dimexpr\linewidth-2\fboxsep}{\small #1}}\par}
\newcommand{\flowbox}[1]{\fcolorbox{ai}{ai!5}{\parbox[c][1.6cm][c]{2.8cm}{\centering\small #1}}}
\newcommand{\bignum}[2]{\parbox[t]{0.235\linewidth}{\centering{\fontsize{22}{26}\selectfont\heiti\color{ai}#1}\\[3pt]{\small #2}}}
\newcommand{\bigval}[1]{{\fontsize{24}{28}\selectfont #1}}

\begin{document}
\begin{titlepage}
\centering\vspace*{45mm}
{\Huge\heiti C3 FPS AI 声卡脚步增强功能\par}\vspace{8pt}
{\huge\heiti 客观评测报告\par}
\vspace{10mm}
{\large 基于 VALORANT 与 PUBG 真实游戏录音的同源三轨对比测试\par}
\vspace{5mm}
{\normalsize 对外简版 · Executive Report\par}
\vspace{35mm}
\begin{tabular}{rl}
评测对象： & Fosi Audio C3 游戏声卡（FPS AI 模式）\\
目标游戏： & VALORANT、PUBG（本阶段计划评测游戏）\\
数据： & 31 段真实对局录音，共 7.7 min\\
编制： & [编制部门 / 人员]\\
版本： & v1.0\quad 2026 年 10 月\\
代码仓库： & \url{https://github.com/21377241/fosi-c3-test}\\
\end{tabular}
\vfill
{\small 本报告为《客观评测技术报告（详细版）》的压缩版本，所有数字与结论均以详细版为准。\par}
\end{titlepage}

\section{结论速览}
\keybox{\textbf{一句话结论：}C3 实际 EQ 在本次测试中并未增强脚步，而 FPS AI 对人工确认脚步的中位增强达到 $+8.7$ dB，同时枪声中位变化仅 $+0.2$ dB。}

\begin{center}
\renewcommand{\arraystretch}{1.6}
\begin{tabular}{lcc}
\toprule
模式 & 脚步中位增益 & 枪声中位增益 \\
\midrule
原始 & \bigval{\color{black!45}0 dB} & \bigval{\color{black!45}0 dB} \\
C3 实际 EQ & \bigval{$-0.8$ dB} & \bigval{$-1.0$ dB} \\
\textbf{\color{ai}C3 FPS AI} & \bigval{\color{ai}\textbf{$+8.7$ dB}} & \bigval{\color{ai}$+0.2$ dB} \\
\bottomrule
\end{tabular}

\vspace{4mm}
\tagbox{脚步更突出}\quad\tagbox{枪声基本保持}\quad\tagbox{固定 EQ 难以复制这种选择性}
\end{center}

两种都是用户在 C3 上实际可用的模式。脚步增强幅度：中位 $+8.7$ dB，94\% $\geq$ 3 dB，整体位于宣传区间（3--9 dB）上沿，部分样本超过 9 dB。

为进一步检验“普通固定 EQ 是否只要调大一点也能做到类似效果”，技术报告额外构造了一个匹配静态 EQ（并非 C3 自带模式），使其脚步增益达到与 FPS AI 相同的 $+8.7$ dB；
此时枪声也被同步提升约 $+8.8$ dB（见第~\ref{sec:core} 节）。

\vspace{3mm}
\greybox{\textbf{边界说明：}不含脚步的非脚步声音中，约 6\% 被误增强超过 3 dB（安静的换弹声、远处枪声），低于实际 EQ 的 17\%\footnotemark；
与脚步重叠的声音会随脚步一起被增强，属于正常；与响亮枪声重叠时枪声基本不受影响。本测试只验证声学处理效果，不直接证明定位能力提升。}
\footnotetext{是否含脚步由只听原始轨的盲审判定。另有 1 例吃药声被增强，据品牌方说明为 FPS AI 的特化增强，不计为误增强。详见技术报告。}

\clearpage
\section{实验为什么公平}
同一段游戏声音同时输出为三条轨，三条轨内容完全相同，差别只来自处理方式。因此可以对\textbf{同一个声音、同一时间点}直接比较“处理后比原始响了多少”。

\begin{center}
\flowbox{\textbf{同一段游戏音频}\\[2pt]VALORANT / PUBG}
$\rightarrow$
\flowbox{\textbf{三轨同步录制}\\[2pt]Raw / FPS AI / EQ}
$\rightarrow$
\flowbox{\textbf{同一个声音}\\[2pt]同一时间点\\比较音量}
$\rightarrow$
\flowbox{\textbf{人工盲听}\\[2pt]确认事件类别}
\end{center}

\begin{itemize}
\item \textbf{Raw}：原始声音，不做处理；
\item \textbf{FPS AI}：开启 C3 的 FPS AI 脚步增强；
\item \textbf{EQ}：C3 产品自带的 EQ 预设（下文称“实际 EQ”）。
\end{itemize}
\vspace{2mm}
\keybox{\textbf{增益怎么算？}对同一个声音事件，分别测原始轨和处理轨在同一时间段内的 RMS（均方根幅值）
\[
\mathrm{RMS}(x)=\sqrt{\frac{1}{N}\sum_{n=1}^{N}x_n^2},\qquad
G=20\log_{10}\frac{\mathrm{RMS}_{\text{处理}}}{\mathrm{RMS}_{\text{原始}}}\ \ (\mathrm{dB})
\]
$G>0$ 表示增强，$G=0$ 表示基本不变；$+6$ dB 约对应波形幅值翻倍。报告中的“脚步 $+8.7$ dB”即人工确认脚步的 $G$ 的中位数。}

\vspace{5mm}
\begin{center}
\bignum{31 段}{真实对局录音}%
\bignum{7.7 min}{总时长}%
\bignum{1620}{自动检出的\\候选声音事件}%
\bignum{348}{人工盲听标注}
\end{center}
\vspace{3mm}

录音来自品牌方提供的 VALORANT 16 段、PUBG 15 段真实对局，全部使用，未按处理效果筛选。主结果仅使用人工确认且无瞬态重叠的事件，包括 170 个脚步和 72 个枪声\footnote{标注采用盲听复核，人机协同模型仅辅助扩展；详细质量控制见技术报告。}。

\clearpage
\section{核心对比：开 C3 实际 EQ 与开 FPS AI，有什么区别？}\label{sec:core}
消费者最直接的问题是：“我直接开 C3 的 EQ，和开 FPS AI，到底有什么区别？”下图比较这两种实际可用的产品模式。

\figslot{figb_core_compare}{0.9\linewidth}{C3 两种实际可用模式对脚步与枪声的影响（人工确认的 170 个脚步、72 个枪声）}{core}

\begin{itemize}
\item \textbf{左：C3 实际 EQ}——脚步 $-0.8$ dB，枪声 $-1.0$ dB：没有增强脚步；
\item \textbf{右：C3 FPS AI}——脚步 $+8.7$ dB，枪声仅 $+0.2$ dB：脚步明显更突出，枪声基本保持。
\end{itemize}
两种模式的作用逻辑完全不同：实际 EQ 是作用于全部声音的固定频响调整，FPS AI 则主要增强脚步。

\figslot{figb_waveforms}{\linewidth}{直观示例：同一个脚步（上）与响亮枪声（下）在三条轨上的波形，同一行纵轴刻度相同}{wave}
FPS AI 轨上脚步波形明显变大（$+9.1$ dB，波形幅值约 ×2.8），枪声几乎不变（$+0.1$ dB）；实际 EQ 轨上两者都没有明显变化（$-1.1$ / $-0.9$ dB）。

\keybox{\textbf{“那如果把 EQ 调得更激进呢？”}\\[3pt]
为回答这个技术追问，技术报告额外构造了一个\textbf{匹配静态 EQ}：它不是 C3 自带模式，而是按 FPS AI 的脚步增强曲线构造、并调到与 FPS AI 相同脚步增益的固定 EQ
（这一设定对 EQ 一方偏有利）。
\begin{center}
\renewcommand{\arraystretch}{1.3}
\begin{tabular}{lcc}
\toprule
 & 脚步 & 枪声 \\
\midrule
C3 FPS AI & $+8.7$ dB & $+0.2$ dB \\
匹配静态 EQ & $+8.7$ dB & \textbf{\color{gun}$+8.8$ dB} \\
\bottomrule
\end{tabular}
\end{center}
即使固定 EQ 强行做到同等脚步增强，枪声也会被一起抬高（背景声 $+8.6$ dB）。为控制“枪声本来比脚步更响”的影响，实验还将每个枪声与原始音量相近（$\pm$3 dB）的脚步进行对照：
FPS AI 下枪声仍比脚步少增强 $6.9$ dB，而匹配静态 EQ 下这一差值仅 $-0.1$ dB。这种差异无法由固定静态频响解释，表现出明显的事件选择性处理特征。}

\section{真实游戏结果：两款目标游戏中均观察到选择性增强}
\begin{table}[htbp]\centering
\caption{两款目标游戏中的中位增益（dB，人工确认事件）}\label{tab:b_game}
{\small\begin{threeparttable}
\begin{tabular}{lcccccc}\toprule
游戏 & \makecell{脚步\\事件数} & \makecell{脚步\\实际 EQ} & \makecell{\textbf{脚步}\\\textbf{FPS AI}} & \makecell{枪声\\事件数} & \makecell{枪声\\实际 EQ} & \makecell{\textbf{枪声}\\\textbf{FPS AI}} \\
\midrule
VALORANT & 97 & $-$1.4 & \hlai{$+$8.0} & 48 & $-$1.4 & \hlai{$+$0.2} \\
PUBG & 73 & $+$0.8 & \hlai{$+$11.8} & 24 & $+$2.7 & \hlai{$+$0.1} \\
\bottomrule\end{tabular}
\begin{tablenotes}\footnotesize\item 加粗：FPS AI 优于同游戏的 C3 实际 EQ（脚步增益更高，枪声更接近 0）。\end{tablenotes}\end{threeparttable}}
\end{table}

两款游戏中实际 EQ 均未明显增强脚步（$-1.4$ / $+0.8$ dB）；FPS AI 的增强幅度不同，但都呈现“脚步明显增强、枪声中位变化较小”的总体模式。

\section{更多证据与已知边界}
\subsection{不仅是音量：FPS AI 对枪声音色结构的改变也明显低于实际 EQ}
\begin{center}
{\fontsize{22}{26}\selectfont\heiti\color{ai}枪声频谱变化 $\approx$ 实际 EQ 的 1/9}\\[4pt]
{\small 频谱总变化：FPS AI $0.9$ dB$^2$ vs 实际 EQ $8.3$ dB$^2$\\
扣除整体音量、只看频谱形状：约 1/6（$0.7$ vs $4.4$ dB$^2$）}
\end{center}
在含响亮枪声的片段中，整段最大音量中位只变化 $+0.1$ dB：最响的声音没有被一起推高。

\figslot{figb_gain_bars}{0.92\linewidth}{各类声音在三种处理下的中位增益（浅蓝带为宣传的 3--9 dB）}{bars}
FPS AI 主要增强脚步，对响亮枪声影响较小；道具类声音（烟雾弹、吃药、爆能器）仅 6 个，且其中 4 个与脚步重叠，其较高的中位增益主要来自重叠的脚步，样本过少，不作结论。

\subsection{已知边界}
\greybox{\begin{itemize}
\item 不含脚步的非脚步声音中，约 6\% 被误增强超过 3 dB，为安静的换弹声和远处枪声，它们与脚步一样安静、短促；
\item 脚步与安静的枪声、道具等声音重叠时，重叠的声音会随脚步一起被增强；与响亮枪声重叠时，枪声基本不受影响；
\item 脚步附近的背景声会随脚步同步抬升；以脚步为主的安静片段，整段最大音量中位上升 $+10.9$ dB；
\item 本测试只验证声学处理效果，不直接证明玩家定位能力提升；
\item 结论限于 VALORANT 与 PUBG 两款目标游戏。
\end{itemize}}

\section{模型更新补充：环境音误增强（PUBG）}
主报告之后，品牌方提供了更新模型。针对官方“优化背景声误增强”口径，在同源 PUBG 素材上对照新旧模型
（仅 PUBG；安静帧口径，与技术报告前文「背景段误增强 30\%」不是同一指标）：
底噪误增强率由 17\% 降至 0\%，中位增益由 $-0.1$ 降至 $-3.5$ dB；
脚步中位由 $+11.8$ 变为 $+12.3$ dB，未见明显削弱。详见技术报告「模型更新：环境音误增强改善」一节及独立短报告。

\section{结论}
\keybox{本次测试表明，在 VALORANT 与 PUBG 两款目标游戏中，C3 实际 EQ 未增强脚步，而 C3 FPS AI 表现出明显的脚步选择性增强：脚步获得显著增益，响亮枪声整体保持稳定；
进一步的技术对照表明，这一差异无法由具有相同脚步增益的固定静态 EQ 复现。}
\vspace{3mm}

实验支持的是声学处理效果。玩家是否因此更早发现敌人、提升定位准确率，仍需进一步主观听辨实验验证。

\vspace{6mm}
{\small 方法细节、标注质量控制、统计检验与完整数据见《C3 FPS AI 声卡脚步增强功能 客观评测技术报告（详细版）》。\\
分析代码与复现脚本：\url{https://github.com/21377241/fosi-c3-test}。\par}
\end{document}
